
\section{Introduction}

\label{sec:intro}

Large Language Models (LLMs) have demonstrated remarkable capabilities across a range of complex tasks~\citep{brown2020language,chowdhery2022palm,touvron2023llama,openai2024gpt4}. This success is largely built on a static ``train then deploy" paradigm, where a model first acquires knowledge from massive corpora and then kept fixed during inference. Yet this design imposes a fundamental limitation: once deployed, the model’s weights cannot be updated, preventing dynamic adaptation to the specific context provided by streaming input tokens. Consequently, at test time, the model is constrained in its ability to process and reason over long-horizon, evolving tasks~\citep{chan2024mle,starace2025paperbench}, and to continuously learn from unbounded streams of experience like humans~\citep{silver2025welcome}.

In-context learning~\citep{brown2020language,wei2023chainofthoughtpromptingelicitsreasoning} offers a way to mitigate this problem via maintaining all past input tokens in the context. However, its effectiveness is tethered to the model's context window, restricted by the quadratic complexity of the de facto attention mechanism~\citep{vaswani2017attention}. This bottleneck has spurred a line of research into architectural solutions aimed at efficiently extending the context window~\citep{beltagy2020longformer,peng2023yarn,child2019generating,dao2023mamba}. Differently, \textit{Test-Time Training (TTT)} has emerged as a new paradigm~\citep{sun2020ttt,wang2021tent,sun2024tttLayers,behrouz2024titans,yau2025sequentialparalleldualityprefixscannable}. Instead of merely making a static model more efficient, TTT enables the model to dynamically update the parameters and adapt to any specific context, directly targeting the aforementioned limitation. Specifically, TTT introduces a small subset of model parameters, called fast weights~\citep{schlag2021linear}, which can be updated on the fly for each new input. By minimizing a self-supervised reconstruction objective, these fast weights compress and internalize contextual information, functioning as an expressive, online evolving state. 

Despite its conceptual appeal, unleashing TTT's potential within the current LLM ecosystem is hindered by critical barriers: \textcolor{myblue}{(\romannumeral1)} Existing TTT methods often rely on specialized layers beyond standard Transformer blocks, which usually demand costly pretraining from scratch to achieve satisfactory performance.~\citep{sun2020ttt,wang2021tent,zhang2025tttdr,sun2024tttLayers}; \textcolor{myblue}{(\romannumeral2)} the canonical TTT mechanism is inherently sequential~\citep{sun2020ttt,sun2024tttLayers}. While existing works explore chunk-wise acceleration~\citep{sun2023retentive,behrouz2024titans,irie2025fastweightprogramminglinear,yau2025sequentialparalleldualityprefixscannable}, TTT's role as the primary token mixer forces a reliance on small chunks to maintain performance, thereby bottlenecking the massive parallelism required to saturate modern accelerators; and \textcolor{myblue}{(\romannumeral3)} the prevalent use of a generic reconstruction objective for TTT's fast weights updating is not explicitly tailored for the causal, Next-Token Prediction task that governs autoregressive LMs, potentially hindering their ultimate performance.

To bridge this gap, we introduce \textit{In-Place Test-Time Training (In-Place TTT)}, a framework designed to seamlessly endow LLMs with Test-Time Training capabilities by directly addressing the aforementioned barriers. Our core insight is to repurpose existing MLP blocks with an in-place design rather than introducing a new, specialized layer \textcolor{myblue}{(tackling barrier \romannumeral1)}. Specifically, In-Place TTT treats the final projection matrix of MLP blocks as their fast weights, updating it in-place during inference. This ``drop-in" design requires no modifications to the model's architecture, preserving the integrity of pre-trained weights and enabling on-the-fly adaptation without costly retraining from scratch.

To tackle the computational inefficiency and objective misalignment, we further design a bespoke adaptation mechanism for language modeling. Following previous works~\citep{sun2023retentive,behrouz2024titans,irie2025fastweightprogramminglinear,zhang2025tttdr,yau2025sequentialparalleldualityprefixscannable}, we replace the inefficient per-token updates with a scalable chunk-wise update rule \textcolor{myblue}{(tackling barrier \romannumeral2)}. Furthermore, our in-place design operates complementarily to the attention mechanism. This synergy obviates the need for small chunks required by standalone TTT layers, thereby ensuring high throughput on modern accelerators. Concurrently, we move beyond the generic reconstruction targets of prior work~\citep{sun2024tttLayers,zhang2025tttdr} and introduce a novel objective explicitly aligned with the Next-Token Prediction (NTP) goal \textcolor{myblue}{(tackling barrier \romannumeral3)}. Grounded in a rigorous theoretical analysis, we show this NTP-aligned objective encourages the fast weights to store predictively useful information for autoregressive language modeling, leading to a highly effective and scalable algorithm.

Grounded in these principled design choices, our In-Place TTT provides a practical and effective framework for enhancing LLMs with dynamic, continual adaptation. We conduct extensive experiments on language modeling tasks of various compute scales, using them as a practical proxy to probe the model's potential on long-horizon, evolving tasks. Through relatively cheap continual training, our In-Place TTT enables Qwen3-4B-Base to achieve superior performance on tasks with contexts up to 128k. Furthermore, we compare In-Place TTT with competitive TTT-related methods by conducting pretraining from scratch on up to 32k-length corpora, validating the architectural merit of our framework. Finally, ablation studies on state size, chunk size, and fast weight objectives provide deeper insights, confirming the critical role of each design choice. Collectively, our results establish In-Place TTT as a promising step towards a paradigm of continual learning in LLMs.